\section{Plan de elicitaci\'on}
\label{sec:plan_elicitacion}
\sloppy

Esta secci\'on describe el plan de sesiones de elicitaci\'on complementaria y de validaci\'on con la organizaci\'on. Los hallazgos ya obtenidos de las cuatro entrevistas de campo se presentan en la secci\'on~6; por ello no se repiten aqu\'i el diagrama de contexto ni la tabla de ocho perspectivas. Las fechas de las sesiones propuestas quedan sujetas a la disponibilidad confirmada por El Balcón Restaurant.

\subsection{Definiciones operativas}
Se entiende por \textbf{parte interesada} a la persona o rol con inter\'es leg\'timo en el sistema a lo largo de su ciclo de vida, incluidos quienes influyen en los objetivos y la operaci\'on \cite{iso29148}. El \textbf{l\'imite del sistema} es la frontera entre los procesos que el SGR-IA abarcar\'a --pedidos, cocina, inventario y ventas-- y las entidades externas que quedan fuera, como la facturaci\'on electr\'onica (secci\'on~\ref{sec:alcance}). El \textbf{contexto} es el entorno de uso en el que esas partes interact\'an, representado en la figura~\ref{fig:contexto}. Una \textbf{t\'ecnica de elicitaci\'on} es el procedimiento para obtener conocimiento de una fuente: seg\'un el syllabus IREB \cite{ireb2026}, la entrevista accede al conocimiento expl\'icito que la persona puede relatar y la observaci\'on revela el conocimiento t\'acito que solo aparece al ejecutar la tarea. Estas definiciones delimitan las decisiones de esta entrega y evitan afirmaciones aplicables a cualquier sistema.

\begin{table}[htbp]
\centering
\footnotesize
\begin{tabularx}{\linewidth}{|X|p{2.5cm}|p{2.8cm}|X|}
\hline
\thdr{Parte interesada} & \thdr{Conocimiento} & \thdr{T\'ecnica} & \thdr{Instrumento}\\
\hline
AD-01 (responsable) & Expl\'icito & Entrevista semiestructurada & Guion de entrevista (A)\\
\hline
COC-01, COC-02 (cocina) & T\'acito y expl\'icito & Observaci\'on y taller & Gu\'ias C y A\\
\hline
CAJ-01 (caja) & Expl\'icito & Entrevista y cuestionario & Guion (A) y cuestionario (B)\\
\hline
INV-01 (inventario) & Expl\'icito & Cuestionario estandarizado & Cuestionario (B)\\
\hline
\end{tabularx}
\caption{Correspondencia entre parte interesada, conocimiento, t\'ecnica e instrumento.}
\label{tab:interesado_tecnica}
\end{table}

La tabla~\ref{tab:interesado_tecnica} resume las decisiones para los roles clave; la tabla~\ref{tab:tecnicas} las extiende al resto de las fuentes identificadas.

\subsection{Selecci\'on de t\'ecnicas seg\'un fuente y tipo de conocimiento}
La selecci\'on considera si el conocimiento buscado es expl\'icito --el que la persona puede relatar-- o t\'acito --el que solo aparece al observar la tarea--, seg\'un los fundamentos del syllabus IREB \cite{ireb2026} y de la ingenier\'ia de requisitos \cite{pohl2025}. La tabla~\ref{tab:tecnicas} vincula cada fuente con la t\'ecnica propuesta.

\begin{table}[htbp]
\centering
\footnotesize
\begin{tabularx}{\linewidth}{|p{2.2cm}|p{3.0cm}|p{2.3cm}|X|}
\hline
\thdr{Fuente (c\'odigo)} & \thdr{T\'ecnica} & \thdr{Conocimiento} & \thdr{Justificaci\'on}\\
\hline
AD-01 & Entrevista semiestructurada & Expl\'icito & Permite obtener objetivos, prioridades y criterios de \'exito directamente del responsable.\\
\hline
GO-01 & Taller de modelado de flujos & T\'acito y expl\'icito & Facilita reconstruir el recorrido del pedido con quien coordina turnos.\\
\hline
COC-01 & Observaci\'on + entrevista corta & T\'acito & Revela esperas y reprocesos que no suelen mencionarse en una entrevista.\\
\hline
COC-02 & Observaci\'on an\'onima y preguntas en taller & T\'acito & Registra reprocesos sin exponer a quien ejecuta la tarea.\\
\hline
MES-01 & Observaci\'on en sala & T\'acito & Captura la transmisi\'on verbal del pedido y sus puntos de p\'erdida.\\
\hline
CAJ-01 & Entrevista focalizada & Expl\'icito & Aporta reglas de cierre, medios de pago y formatos de resumen.\\
\hline
INV-01 & Cuestionario & Expl\'icito & Recoge entradas, salidas y frecuencia de conteo de forma estandarizada.\\
\hline
DEL-01 & Cuestionario breve & Expl\'icito & Describe el estado de entrega sin exigir una sesi\'on presencial.\\
\hline
CLI-01 & Cuestionario an\'onimo & Expectativas & Obtiene expectativas de servicio sin exponer datos personales.\\
\hline
PRO-01 & Contacto indirecto por INV-01 & Expl\'icito & Recoge plazos de reposici\'on sin solicitar datos del proveedor.\\
\hline
\end{tabularx}
\caption{T\'ecnicas de elicitaci\'on por fuente y tipo de conocimiento.}
\label{tab:tecnicas}
\end{table}

\subsection{Cronograma de sesiones propuesto}
Las fechas de la tabla~\ref{tab:sesiones} son una propuesta de las semanas 5 y 6 del SPA 2026--2027 y quedan sujetas a la disponibilidad confirmada por El Balcón Restaurant.

\begin{table}[htbp]
\centering
\footnotesize
\begin{tabularx}{\linewidth}{|l|p{1.9cm}|p{2.6cm}|p{2.7cm}|X|}
\hline
\thdr{Sesi\'on} & \thdr{Semana} & \thdr{Participantes} & \thdr{T\'ecnica} & \thdr{Producto esperado}\\
\hline
S1 & 5 & AD-01, GO-01 & Entrevista y taller & Visi\'on y flujo actual acordados.\\
\hline
S2 & 5 & MES-01, COC-01 & Observaci\'on en turno & Notas an\'onimas de tiempos y reprocesos.\\
\hline
S3 & 6 & CAJ-01, INV-01 & Entrevista y cuestionario & Reglas de caja e inventario.\\
\hline
S4 & 6 & DEL-01, CLI-01 & Cuestionario breve & Expectativas de entrega y servicio.\\
\hline
S5 & 6 & AD-01, GO-01 & Validaci\'on & Revisiones al contexto y a las ocho perspectivas.\\
\hline
\end{tabularx}
\caption{Cronograma propuesto de sesiones de elicitaci\'on.}
\label{tab:sesiones}
\end{table}

\subsection{Instrumentos de apoyo}
Se entregan tres instrumentos en los ap\'endices: el guion de entrevista (Ap\'endice~A), el cuestionario (Ap\'endice~B) y la gu\'ia de observaci\'on y taller (Ap\'endice~C). Los tres usan c\'odigos de rol, recogen exclusivamente lo necesario para la elicitaci\'on y explican al participante el tratamiento de sus datos antes de aplicarse, de acuerdo con la LOPDP \cite{lopdp2021}.
